\section{Diagonalization of symmetric matrices}
A real matrix is symmetric when $A^T=A$, and an orthogonal matrix satisfies
$P^TP=I$, so $P^{-1}=P^T$. These are different definitions.
The spectral theorem says that a real square matrix has an orthonormal
eigenvector basis exactly when it is symmetric.
To orthogonally diagonalize a symmetric matrix:
\begin{enumerate}
\item Find its real eigenvalues, with multiplicities.
\item For each $\lambda$, solve $(A-\lambda I)v=0$ by row reduction.
\item Within a repeated eigenspace, use Gram--Schmidt if needed.
Eigenvectors belonging to distinct eigenvalues are already perpendicular.
\item Normalize each nonzero vector $v_i$ to $u_i=v_i/\|v_i\|$.
Set $P=[u_1\ \cdots\ u_n]$ and $D=\operatorname{diag}(\lambda_1,\ldots,\lambda_n)$
in the same order. Then $A=PDP^T$.
\end{enumerate}
In the computations below, the columns of $V$ are unnormalized orthogonal
eigenvectors. A diagonal factor rescales the columns to form $P$.
Every displayed diagonalization was checked exactly by $AP=PD$ and $P^TP=I$.
Changing a column's sign or reordering columns together with $D$ gives
another equally valid answer.

\practice{7.1.P1}{A square of a symmetric matrix is symmetric}
Using the reverse-order transpose rule,
$(A^2)^T=(AA)^T=A^TA^T=AA=A^2$, since $A^T=A$.
Thus $A^2$ is symmetric. The reversal creates no difficulty here because
the two factors are the same matrix.

\practice{7.1.P2}{Square an orthogonal diagonalization}
If $A=PDP^T$ with $P^TP=I$, then
$A^2=PDP^TPDP^T=PD^2P^T$.
The same $P$ is orthogonal and $D^2$ is diagonal, so this is an orthogonal
diagonalization of $A^2$.

\exercise{7.1.1}{Match entries across the diagonal}
The two off-diagonal entries are both $5$, so
$\begin{bmatrix}3&5\\5&-7\end{bmatrix}^T$
is the original matrix. It is symmetric.

\exercise{7.1.2}{Opposite signs do not match}
The upper-right entry is $5$ and the lower-left entry is $-5$.
They would have to be equal for symmetry. Therefore the matrix is not symmetric.

\exercise{7.1.3}{One unequal pair is enough}
Here $a_{12}=2$ but $a_{21}=4$. Thus $A^T\ne A$, so the matrix is
not symmetric, despite the repeated entries within each row.

\exercise{7.1.4}{Check all three off-diagonal pairs}
The paired entries are $a_{12}=a_{21}=8$,
$a_{13}=a_{31}=3$, and $a_{23}=a_{32}=-2$.
The diagonal is unchanged by transpose. Thus this matrix is symmetric.

\exercise{7.1.5}{An upper triangular matrix with nonzero off-diagonal entries}
The entry $a_{12}=2$ has counterpart $a_{21}=0$.
Thus the matrix is not symmetric. An upper triangular symmetric matrix
would have to be diagonal, because its entries below and above the diagonal
must both vanish.

\exercise{7.1.6}{Symmetry requires a square matrix}
The given array has three rows and four columns. Its transpose has four
rows and three columns, so the two cannot be equal. It is not symmetric.
The repeated numerical pattern does not overcome the shape mismatch.

\exercise{7.1.7}{Orthonormal decimal columns}
The two columns have squared lengths $.6^2+.8^2=1$ and
$.8^2+(-.6)^2=1$. Their dot product is $.6(.8)+.8(-.6)=0$.
Thus $A$ is orthogonal, and
\[
 A^{-1}=A^T=\begin{bmatrix}.6&.8\\.8&-.6\end{bmatrix}.
\]
Here the matrix happens to be symmetric too, so it is its own inverse.

\exercise{7.1.8}{A rotation has its transpose as inverse}
Each column has squared length $1/2+1/2=1$, and the dot product is
$-1/2+1/2=0$. Hence the matrix is orthogonal and
\[
 A^{-1}=A^T=\frac1{\sqrt2}\begin{bmatrix}1&1\\-1&1\end{bmatrix}.
\]
Geometrically this reverses the $45^\circ$ rotation of the original matrix.

\exercise{7.1.9}{Perpendicular but not unit columns}
The columns $(-5,2)^T$ and $(2,5)^T$ are perpendicular, but each has
squared length $25+4=29$. Thus $A^TA=29I$, not $I$.
The matrix is not orthogonal. Having an inverse is a separate issue;
orthogonality specifically requires the columns to be normalized.

\exercise{7.1.10}{Another missing normalization}
Each column has squared length $1+4+4=9$, and distinct columns have
dot product $-2-2+4=0$. Therefore $A^TA=9I$, not $I$.
The given matrix is not orthogonal; dividing it by $3$ would make it so.

\exercise{7.1.11}{Verify a less obvious orthogonal matrix}
Let the columns be $c_1,c_2,c_3$. Their squared lengths are
\[
 \|c_1\|^2=4/9+5/9=1,\quad
 \|c_2\|^2=4/9+1/5+16/45=1,\quad
 \|c_3\|^2=1/9+4/5+4/45=1.
\]
Their pairwise dot products are
$4/9-4/9=0$, $2/9-2/9=0$, and
$2/9-2/5+8/45=0$.
Thus all columns are orthonormal and the inverse is the transpose:
\[
 A^{-1}=\begin{bmatrix}
 2/3&0&\sqrt5/3\\2/3&1/\sqrt5&-4/\sqrt{45}\\
 1/3&-2/\sqrt5&-2/\sqrt{45}
 \end{bmatrix}.
\]

\exercise{7.1.12}{An orthogonal signed matrix}
Write $A=H/2$, where $H$ has the printed entries $1$ or $-1$.
Each column of $H$ has squared length $4$, and every pair of different
columns has two agreeing signs and two disagreeing signs, so dot product zero.
Thus $A^TA=I$, and
\[
 A^{-1}=A^T=\frac12\begin{bmatrix}
 1&-1&1&-1\\1&1&1&1\\-1&-1&1&1\\-1&1&1&-1
 \end{bmatrix}.
\]

\exercise{7.1.13}{An orthogonal eigenvector basis}
The characteristic polynomial $\det(A-tI)$ is $\left(t - 4\right) \left(t - 2\right)$. Its roots give the eigenvalues below.
\step{Solve the eigenspace equations.}
$\lambda=2$: row reduction gives \(x_{1}=- x_{2}\).
$\lambda=4$: row reduction gives \(x_{1}=x_{2}\).
\step{Normalize the orthogonal vectors and pair the same order with D.}
The columns of
\[
 V=\begin{bmatrix}1&1\\-1&1\end{bmatrix}
\]
are mutually orthogonal eigenvectors in the order
\[
 (\lambda_1,\ldots,\lambda_2)=(2,4).
\]
Their squared lengths are $(2,2)$. Hence an explicit answer is
\[
 P=\frac1{\sqrt{2}}V,\qquad
 D=\operatorname{diag}(2,4).
\]
Dividing columns by these nonzero lengths gives $P^TP=I$, and the
eigenspace equations give $AP=PD$. Multiplying the latter by $P^T$
therefore verifies $A=PDP^T$.

\exercise{7.1.14}{An orthogonal eigenvector basis}
The characteristic polynomial $\det(A-tI)$ is $\left(t - 6\right) \left(t + 4\right)$. Its roots give the eigenvalues below.
\step{Solve the eigenspace equations.}
$\lambda=-4$: row reduction gives \(x_{1}=- x_{2}\).
$\lambda=6$: row reduction gives \(x_{1}=x_{2}\).
\step{Normalize the orthogonal vectors and pair the same order with D.}
The columns of
\[
 V=\begin{bmatrix}1&1\\-1&1\end{bmatrix}
\]
are mutually orthogonal eigenvectors in the order
\[
 (\lambda_1,\ldots,\lambda_2)=(-4,6).
\]
Their squared lengths are $(2,2)$. Hence an explicit answer is
\[
 P=\frac1{\sqrt{2}}V,\qquad
 D=\operatorname{diag}(-4,6).
\]
Dividing columns by these nonzero lengths gives $P^TP=I$, and the
eigenspace equations give $AP=PD$. Multiplying the latter by $P^T$
therefore verifies $A=PDP^T$.

\exercise{7.1.15}{An orthogonal eigenvector basis}
The characteristic polynomial $\det(A-tI)$ is $t \left(t - 17\right)$. Its roots give the eigenvalues below.
\step{Solve the eigenspace equations.}
$\lambda=0$: row reduction gives \(x_{1}=\frac{x_{2}}{4}\).
$\lambda=17$: row reduction gives \(x_{1}=- 4 x_{2}\).
\step{Normalize the orthogonal vectors and pair the same order with D.}
The columns of
\[
 V=\begin{bmatrix}1&4\\4&-1\end{bmatrix}
\]
are mutually orthogonal eigenvectors in the order
\[
 (\lambda_1,\ldots,\lambda_2)=(0,17).
\]
Their squared lengths are $(17,17)$. Hence an explicit answer is
\[
 P=\frac1{\sqrt{17}}V,\qquad
 D=\operatorname{diag}(0,17).
\]
Dividing columns by these nonzero lengths gives $P^TP=I$, and the
eigenspace equations give $AP=PD$. Multiplying the latter by $P^T$
therefore verifies $A=PDP^T$.

\exercise{7.1.16}{An orthogonal eigenvector basis}
The characteristic polynomial $\det(A-tI)$ is $\left(t - 25\right) \left(t + 25\right)$. Its roots give the eigenvalues below.
\step{Solve the eigenspace equations.}
$\lambda=-25$: row reduction gives \(x_{1}=- \frac{4 x_{2}}{3}\).
$\lambda=25$: row reduction gives \(x_{1}=\frac{3 x_{2}}{4}\).
\step{Normalize the orthogonal vectors and pair the same order with D.}
The columns of
\[
 V=\begin{bmatrix}4&3\\-3&4\end{bmatrix}
\]
are mutually orthogonal eigenvectors in the order
\[
 (\lambda_1,\ldots,\lambda_2)=(-25,25).
\]
Their squared lengths are $(25,25)$. Hence an explicit answer is
\[
 P=\frac1{5}V,\qquad
 D=\operatorname{diag}(-25,25).
\]
Dividing columns by these nonzero lengths gives $P^TP=I$, and the
eigenspace equations give $AP=PD$. Multiplying the latter by $P^T$
therefore verifies $A=PDP^T$.

\exercise{7.1.17}{An orthogonal eigenvector basis}
\step{Solve the eigenspace equations.}
$\lambda=-2$: row reduction gives \(x_{1}=- x_{3},\quad x_{2}=0\).
$\lambda=2$: row reduction gives \(x_{1}=x_{3},\quad x_{2}=- 2 x_{3}\).
$\lambda=5$: row reduction gives \(x_{1}=x_{3},\quad x_{2}=x_{3}\).
\step{Normalize the orthogonal vectors and pair the same order with D.}
The columns of
\[
 V=\begin{bmatrix}1&1&1\\0&-2&1\\-1&1&1\end{bmatrix}
\]
are mutually orthogonal eigenvectors in the order
\[
 (\lambda_1,\ldots,\lambda_3)=(-2,2,5).
\]
Their squared lengths are $(2,6,3)$. Hence an explicit answer is
\[
 P=V\operatorname{diag}(\frac{\sqrt{2}}{2},\frac{\sqrt{6}}{6},\frac{\sqrt{3}}{3}),\qquad
 D=\operatorname{diag}(-2,2,5).
\]
Dividing columns by these nonzero lengths gives $P^TP=I$, and the
eigenspace equations give $AP=PD$. Multiplying the latter by $P^T$
therefore verifies $A=PDP^T$.

\exercise{7.1.18}{An orthogonal eigenvector basis}
\step{Solve the eigenspace equations.}
$\lambda=-50$: row reduction gives \(x_{1}=\frac{3 x_{2}}{4},\quad x_{3}=0\).
$\lambda=3$: row reduction gives \(x_{1}=0,\quad x_{2}=0\).
$\lambda=25$: row reduction gives \(x_{1}=- \frac{4 x_{2}}{3},\quad x_{3}=0\).
\step{Normalize the orthogonal vectors and pair the same order with D.}
The columns of
\[
 V=\begin{bmatrix}3&0&4\\4&0&-3\\0&1&0\end{bmatrix}
\]
are mutually orthogonal eigenvectors in the order
\[
 (\lambda_1,\ldots,\lambda_3)=(-50,3,25).
\]
Their squared lengths are $(25,1,25)$. Hence an explicit answer is
\[
 P=V\operatorname{diag}(\frac{1}{5},1,\frac{1}{5}),\qquad
 D=\operatorname{diag}(-50,3,25).
\]
Dividing columns by these nonzero lengths gives $P^TP=I$, and the
eigenspace equations give $AP=PD$. Multiplying the latter by $P^T$
therefore verifies $A=PDP^T$.

\exercise{7.1.19}{An orthogonal eigenvector basis}
\step{Solve the eigenspace equations.}
$\lambda=-2$: row reduction gives \(x_{1}=- x_{3},\quad x_{2}=- \frac{x_{3}}{2}\).
$\lambda=7$: row reduction gives \(x_{1}=- \frac{x_{2}}{2} + x_{3}\).
Within the $7$-eigenspace, start with $(1,0,1)^T$ and $(-1/2,1,0)^T$.
Subtracting the projection of the latter onto the former gives
$(-1/4,1,1/4)^T$, a multiple of $(-1,4,1)^T$. This is the Gram--Schmidt
step that produces the two orthogonal columns used below.
\step{Normalize the orthogonal vectors and pair the same order with D.}
The columns of
\[
 V=\begin{bmatrix}2&1&-1\\1&0&4\\-2&1&1\end{bmatrix}
\]
are mutually orthogonal eigenvectors in the order
\[
 (\lambda_1,\ldots,\lambda_3)=(-2,7,7).
\]
Their squared lengths are $(9,2,18)$. Hence an explicit answer is
\[
 P=V\operatorname{diag}(\frac{1}{3},\frac{\sqrt{2}}{2},\frac{\sqrt{2}}{6}),\qquad
 D=\operatorname{diag}(-2,7,7).
\]
Dividing columns by these nonzero lengths gives $P^TP=I$, and the
eigenspace equations give $AP=PD$. Multiplying the latter by $P^T$
therefore verifies $A=PDP^T$.

\exercise{7.1.20}{An orthogonal eigenvector basis}
\step{Solve the eigenspace equations.}
$\lambda=1$: row reduction gives \(x_{1}=- 2 x_{3},\quad x_{2}=- 2 x_{3}\).
$\lambda=7$: row reduction gives \(x_{1}=- \frac{x_{3}}{2},\quad x_{2}=x_{3}\).
$\lambda=13$: row reduction gives \(x_{1}=x_{3},\quad x_{2}=- \frac{x_{3}}{2}\).
\step{Normalize the orthogonal vectors and pair the same order with D.}
The columns of
\[
 V=\begin{bmatrix}2&1&2\\2&-2&-1\\-1&-2&2\end{bmatrix}
\]
are mutually orthogonal eigenvectors in the order
\[
 (\lambda_1,\ldots,\lambda_3)=(1,7,13).
\]
Their squared lengths are $(9,9,9)$. Hence an explicit answer is
\[
 P=\frac1{3}V,\qquad
 D=\operatorname{diag}(1,7,13).
\]
Dividing columns by these nonzero lengths gives $P^TP=I$, and the
eigenspace equations give $AP=PD$. Multiplying the latter by $P^T$
therefore verifies $A=PDP^T$.

\exercise{7.1.21}{An orthogonal eigenvector basis}
\step{Solve the eigenspace equations.}
$\lambda=1$: row reduction gives \(x_{1}=- x_{3},\quad x_{2}=- x_{4}\).
$\lambda=5$: row reduction gives \(x_{1}=- x_{4},\quad x_{2}=x_{4},\quad x_{3}=- x_{4}\).
$\lambda=9$: row reduction gives \(x_{1}=x_{4},\quad x_{2}=x_{4},\quad x_{3}=x_{4}\).
\step{Normalize the orthogonal vectors and pair the same order with D.}
The columns of
\[
 V=\begin{bmatrix}1&0&1&1\\0&1&-1&1\\-1&0&1&1\\0&-1&-1&1\end{bmatrix}
\]
are mutually orthogonal eigenvectors in the order
\[
 (\lambda_1,\ldots,\lambda_4)=(1,1,5,9).
\]
Their squared lengths are $(2,2,4,4)$. Hence an explicit answer is
\[
 P=V\operatorname{diag}(\frac{\sqrt{2}}{2},\frac{\sqrt{2}}{2},\frac{1}{2},\frac{1}{2}),\qquad
 D=\operatorname{diag}(1,1,5,9).
\]
Dividing columns by these nonzero lengths gives $P^TP=I$, and the
eigenspace equations give $AP=PD$. Multiplying the latter by $P^T$
therefore verifies $A=PDP^T$.

\exercise{7.1.22}{An orthogonal eigenvector basis}
\step{Solve the eigenspace equations.}
$\lambda=0$: row reduction gives \(x_{1}=0,\quad x_{2}=- x_{4},\quad x_{3}=0\).
$\lambda=2$: row reduction gives \(x_{2}=x_{4}\).
\step{Normalize the orthogonal vectors and pair the same order with D.}
The columns of
\[
 V=\begin{bmatrix}0&1&0&0\\1&0&0&1\\0&0&1&0\\-1&0&0&1\end{bmatrix}
\]
are mutually orthogonal eigenvectors in the order
\[
 (\lambda_1,\ldots,\lambda_4)=(0,2,2,2).
\]
Their squared lengths are $(2,1,1,2)$. Hence an explicit answer is
\[
 P=V\operatorname{diag}(\frac{\sqrt{2}}{2},1,1,\frac{\sqrt{2}}{2}),\qquad
 D=\operatorname{diag}(0,2,2,2).
\]
Dividing columns by these nonzero lengths gives $P^TP=I$, and the
eigenspace equations give $AP=PD$. Multiplying the latter by $P^T$
therefore verifies $A=PDP^T$.

\exercise{7.1.23}{An orthogonal eigenvector basis}
The printed vector $v=(1,1,1)^T$ satisfies $Av=5v$, so its eigenvalue
is $5$. The exercise also asks to verify $2$ is an eigenvalue; the vector
$(1,-1,0)^T$ satisfies $A(1,-1,0)^T=2(1,-1,0)^T$.
Thus the two assertions in the wording refer to different eigenvectors.
\step{Solve the eigenspace equations.}
$\lambda=2$: row reduction gives \(x_{1}=- x_{2} - x_{3}\).
$\lambda=5$: row reduction gives \(x_{1}=x_{3},\quad x_{2}=x_{3}\).
\step{Normalize the orthogonal vectors and pair the same order with D.}
The columns of
\[
 V=\begin{bmatrix}1&1&1\\-1&1&1\\0&-2&1\end{bmatrix}
\]
are mutually orthogonal eigenvectors in the order
\[
 (\lambda_1,\ldots,\lambda_3)=(2,2,5).
\]
Their squared lengths are $(2,6,3)$. Hence an explicit answer is
\[
 P=V\operatorname{diag}(\frac{\sqrt{2}}{2},\frac{\sqrt{6}}{6},\frac{\sqrt{3}}{3}),\qquad
 D=\operatorname{diag}(2,2,5).
\]
Dividing columns by these nonzero lengths gives $P^TP=I$, and the
eigenspace equations give $AP=PD$. Multiplying the latter by $P^T$
therefore verifies $A=PDP^T$.

\exercise{7.1.24}{An orthogonal eigenvector basis}
Multiplication gives $Av_1=10v_1$ for the supplied $v_1=(-2,2,1)^T$,
and $Av_2=v_2$ for $v_2=(1,1,0)^T$. For the repeated eigenvalue $1$,
$(1,-1,4)^T$ is a second eigenvector orthogonal to $v_2$.
\step{Solve the eigenspace equations.}
$\lambda=1$: row reduction gives \(x_{1}=x_{2} + \frac{x_{3}}{2}\).
$\lambda=10$: row reduction gives \(x_{1}=- 2 x_{3},\quad x_{2}=2 x_{3}\).
\step{Normalize the orthogonal vectors and pair the same order with D.}
The columns of
\[
 V=\begin{bmatrix}1&1&2\\1&-1&-2\\0&4&-1\end{bmatrix}
\]
are mutually orthogonal eigenvectors in the order
\[
 (\lambda_1,\ldots,\lambda_3)=(1,1,10).
\]
Their squared lengths are $(2,18,9)$. Hence an explicit answer is
\[
 P=V\operatorname{diag}(\frac{\sqrt{2}}{2},\frac{\sqrt{2}}{6},\frac{1}{3}),\qquad
 D=\operatorname{diag}(1,1,10).
\]
Dividing columns by these nonzero lengths gives $P^TP=I$, and the
eigenspace equations give $AP=PD$. Multiplying the latter by $P^T$
therefore verifies $A=PDP^T$.

\exercise{7.1.25}{Symmetry, multiplicity, and a missing unit condition}
\begin{enumerate}[label=(\alph*)]
\item \textbf{True.} If $A=PDP^T$, then
$A^T=(PDP^T)^T=PD^TP^T=PDP^T=A$.
\item \textbf{True.} Symmetry gives
$3u^Tv=(Au)^Tv=u^TAv=4u^Tv$, so $u^Tv=0$.
\item \textbf{False.} $I_n$ is symmetric but has only one distinct
eigenvalue, namely $1$ with multiplicity $n$.
\item \textbf{False.} For $B=vv^T$, $B^2=(v^Tv)B$.
A projection requires $B^2=B$, which for nonzero $v$ requires $\|v\|=1$.
For $v=2e_1$, $B=4e_1e_1^T$ is not idempotent.
\end{enumerate}

\exercise{7.1.26}{Four consequences and nonconsequences of the spectral theorem}
\begin{enumerate}[label=(\alph*)]
\item \textbf{True.} This is the real symmetric spectral theorem.
\item \textbf{True.} Taking the transpose of $PDP^T$ gives the same
matrix because $D$ is real diagonal.
\item \textbf{False.} A $90^\circ$ rotation
$\begin{bmatrix}0&-1\\1&0\end{bmatrix}$ is orthogonal but not symmetric,
so it cannot be orthogonally diagonalized over $\R$.
\item \textbf{True.} In the spectral theorem the orthogonal
diagonalization provides one basis vector for each occurrence of an
eigenvalue on the diagonal, giving equality of geometric and algebraic multiplicities.
\end{enumerate}

\exercise{7.1.27}{Three symmetric products}
Use the reverse-order transpose rule and the hypothesis $A^T=A$:
\[
 (B^TAB)^T=B^TA^TB=B^TAB,
 \quad (B^TB)^T=B^TB,
 \quad (BB^T)^T=BB^T.
\]
The products have square sizes $m\times m$, $m\times m$, and $n\times n$,
respectively, so these equalities indeed establish symmetry in each case.

\exercise{7.1.28}{Move a symmetric matrix across a dot product}
For every $x,y$,
$(Ax)\cdot y=(Ax)^Ty=x^TA^Ty=x^TAy=x\cdot(Ay)$.
The step using $A^T=A$ is exactly where symmetry matters.
Without it the correct expression would contain $A^T$ on the other side.

\exercise{7.1.29}{The inverse retains the eigenvectors}
Write $A=PDP^T$ with $P$ orthogonal. Since $A$ is invertible, none
of the diagonal eigenvalues in $D$ is zero, so $D^{-1}$ exists and is diagonal.
Then $A^{-1}=PD^{-1}P^T$, an orthogonal diagonalization with the same
$P$. Its diagonal eigenvalues are the reciprocals of the original ones.

\exercise{7.1.30}{A commuting product remains symmetric}
Orthogonal diagonalizability makes both $A$ and $B$ symmetric.
Therefore $(AB)^T=B^TA^T=BA=AB$, using the additional commutation hypothesis.
Thus $AB$ is symmetric and the spectral theorem makes it orthogonally
diagonalizable. Without $AB=BA$, a product of two symmetric matrices need
not be symmetric.

\exercise{7.1.31}{Count a repeated eigenvalue's eigenspace dimension}
Let $x=Py$. Since $P$ is invertible,
$Ax=\lambda x$ is equivalent to $(D-\lambda I)y=0$.
At positions where the diagonal of $D$ differs from $\lambda$, this forces
$y_i=0$. At its $k$ positions equal to $\lambda$, the corresponding
$y_i$ are arbitrary. Thus the solution space has exactly $k$ independent
coordinates. Mapping by $P$ gives the eigenspace of $A$ and preserves
its dimension, so that dimension is $k$.
Equivalently the $k$ corresponding columns of $P$ form its basis.

\exercise{7.1.32}{A symmetric triangular matrix is diagonal}
From $A=PRP^T$, solve for $R=P^TAP$.
Taking transposes gives $R^T=P^TA^TP=P^TAP=R$, so $R$ is symmetric.
But $R$ is also upper triangular. Entries below its diagonal are zero;
symmetry forces their counterpart entries above the diagonal to be zero too.
Thus all off-diagonal entries vanish and $R$ is diagonal.

\exercise{7.1.33}{A spectral decomposition for Example 2}
Use the orthogonal eigenvectors
$a=(-1,1,0)^T$, $b=(-1,-1,2)^T$, $c=(1,1,1)^T$,
with eigenvalues $8,6,3$ and squared lengths $2,6,3$.
After normalization the spectral formula gives
\[
 \boxed{A=4aa^T+bb^T+cc^T}.
\]
Indeed the normalization divides each outer product by its squared length;
thus the three coefficients are $8/2$, $6/6$, $3/3$.
Expanding these outer products gives
$\begin{bmatrix}6&-2&-1\\-2&6&-1\\-1&-1&5\end{bmatrix}$,
the matrix from Example 2. This checks the decomposition entrywise.

\exercise{7.1.34}{A spectral decomposition with a repeated eigenvalue}
For Example 3, take $a=(1,0,1)^T$, $b=(-1,4,1)^T$, and
$c=(-2,-1,2)^T$. They are mutually perpendicular, with squared lengths
$2,18,9$. The first two have eigenvalue $7$, and the last has eigenvalue
$-2$. Therefore
\[
 \boxed{A=\frac72aa^T+\frac7{18}bb^T-\frac29cc^T}.
\]
For instance the $(1,1)$ entry is $7/2+7/18-8/9=3$;
expansion of all entries returns
$\begin{bmatrix}3&-2&4\\-2&6&2\\4&2&3\end{bmatrix}$.
The two separate rank-one terms for eigenvalue $7$ reflect its
two-dimensional eigenspace.

\exercise{7.1.35}{A rank-one projection from a unit vector}
\begin{enumerate}[label=(\alph*)]
\item $Bx=(uu^T)x=u(u^Tx)=(x\cdot u)u$.
Since $u\cdot u=1$, this equals
$((x\cdot u)/(u\cdot u))u$, the orthogonal projection formula.
\item $B^T=(uu^T)^T=uu^T=B$ and
$B^2=u(u^Tu)u^T=uu^T=B$.
\item $Bu=u(u^Tu)=u$. The unit vector is nonzero, so it is an
eigenvector with eigenvalue $1$.
\end{enumerate}

\exercise{7.1.36}{Every symmetric idempotent matrix is an orthogonal projection}
\step{(a) Show the proposed pieces are perpendicular.}
Put $\widehat y=By$, $z=(I-B)y$. Symmetry and idempotence give
\[
 z^T\widehat y=y^T(I-B)^TBy=y^T(B-B^2)y=0.
\]
\step{(b) Verify the stronger orthogonality needed for a projection.}
Clearly $\widehat y\in W=\Col B$. For any $w\in W$, write $w=Bv$;
then $z^Tw=y^T(I-B)Bv=0$. Thus $z\in W^\perp$, not merely
perpendicular to one particular vector $\widehat y$.
The identity $y=By+(y-By)$ is therefore the orthogonal decomposition
into a vector in $W$ and one in $W^\perp$. Its uniqueness proves that
$By=\operatorname{proj}_W y$.

\exercise{7.1.37}{Compute eigenvalues, then row-reduce each eigenspace}
The eigenvalues were computed from the characteristic polynomial. For
the eigenvectors we solve null-space systems by row reduction, not by
calling an eigenvector routine, as the exercise requests.
\step{Solve the eigenspace equations.}
$\lambda=-12$: row reduction gives \(x_{1}=x_{4},\quad x_{2}=- x_{4},\quad x_{3}=- x_{4}\).
$\lambda=4$: row reduction gives \(x_{1}=- x_{4},\quad x_{2}=- x_{4},\quad x_{3}=x_{4}\).
$\lambda=10$: row reduction gives \(x_{1}=x_{4},\quad x_{2}=x_{4},\quad x_{3}=x_{4}\).
$\lambda=18$: row reduction gives \(x_{1}=- x_{4},\quad x_{2}=x_{4},\quad x_{3}=- x_{4}\).
\step{Normalize the orthogonal vectors and pair the same order with D.}
The columns of
\[
 V=\begin{bmatrix}1&1&1&1\\-1&1&1&-1\\-1&-1&1&1\\1&-1&1&-1\end{bmatrix}
\]
are mutually orthogonal eigenvectors in the order
\[
 (\lambda_1,\ldots,\lambda_4)=(-12,4,10,18).
\]
Their squared lengths are $(4,4,4,4)$. Hence an explicit answer is
\[
 P=\frac1{2}V,\qquad
 D=\operatorname{diag}(-12,4,10,18).
\]
Dividing columns by these nonzero lengths gives $P^TP=I$, and the
eigenspace equations give $AP=PD$. Multiplying the latter by $P^T$
therefore verifies $A=PDP^T$.

\exercise{7.1.38}{Compute eigenvalues, then row-reduce each eigenspace}
The eigenvalues were computed from the characteristic polynomial. For
the eigenvectors we solve null-space systems by row reduction, not by
calling an eigenvector routine, as the exercise requests.
\step{Solve the eigenspace equations.}
$\lambda=\frac{1}{4}$: row reduction gives \(x_{1}=4 x_{4},\quad x_{2}=2 x_{4},\quad x_{3}=2 x_{4}\).
$\lambda=\frac{3}{10}$: row reduction gives \(x_{1}=- \frac{x_{4}}{4},\quad x_{2}=- \frac{x_{4}}{2},\quad x_{3}=\frac{x_{4}}{2}\).
$\lambda=\frac{11}{20}$: row reduction gives \(x_{1}=x_{4},\quad x_{2}=- \frac{x_{4}}{2},\quad x_{3}=- 2 x_{4}\).
$\lambda=\frac{3}{4}$: row reduction gives \(x_{1}=- x_{4},\quad x_{2}=2 x_{4},\quad x_{3}=- \frac{x_{4}}{2}\).
\step{Normalize the orthogonal vectors and pair the same order with D.}
The columns of
\[
 V=\begin{bmatrix}4&1&2&2\\2&2&-1&-4\\2&-2&-4&1\\1&-4&2&-2\end{bmatrix}
\]
are mutually orthogonal eigenvectors in the order
\[
 (\lambda_1,\ldots,\lambda_4)=(\frac{1}{4},\frac{3}{10},\frac{11}{20},\frac{3}{4}).
\]
Their squared lengths are $(25,25,25,25)$. Hence an explicit answer is
\[
 P=\frac1{5}V,\qquad
 D=\operatorname{diag}(\frac{1}{4},\frac{3}{10},\frac{11}{20},\frac{3}{4}).
\]
Dividing columns by these nonzero lengths gives $P^TP=I$, and the
eigenspace equations give $AP=PD$. Multiplying the latter by $P^T$
therefore verifies $A=PDP^T$.

\exercise{7.1.39}{Compute eigenvalues, then row-reduce each eigenspace}
The eigenvalues were computed from the characteristic polynomial. For
the eigenvectors we solve null-space systems by row reduction, not by
calling an eigenvector routine, as the exercise requests.
\step{Solve the eigenspace equations.}
$\lambda=- \frac{5}{4}$: row reduction gives \(x_{1}=- x_{4},\quad x_{2}=2 x_{4},\quad x_{3}=- \frac{x_{4}}{2}\).
$\lambda=0$: row reduction gives \(x_{1}=- x_{4},\quad x_{2}=- \frac{x_{4}}{2},\quad x_{3}=2 x_{4}\).
$\lambda=\frac{3}{4}$: row reduction gives \(x_{1}=\frac{3 x_{3}}{2} + x_{4},\quad x_{2}=x_{3}\).
For the repeated eigenvalue $3/4$, the two chosen null vectors
$(4,2,2,1)^T$ and $(1,-2,-2,4)^T$ have dot product $4-4-4+4=0$.
Both have squared length $25$, so dividing each by $5$ normalizes them.
\step{Normalize the orthogonal vectors and pair the same order with D.}
The columns of
\[
 V=\begin{bmatrix}2&2&4&1\\-4&1&2&-2\\1&-4&2&-2\\-2&-2&1&4\end{bmatrix}
\]
are mutually orthogonal eigenvectors in the order
\[
 (\lambda_1,\ldots,\lambda_4)=(- \frac{5}{4},0,\frac{3}{4},\frac{3}{4}).
\]
Their squared lengths are $(25,25,25,25)$. Hence an explicit answer is
\[
 P=\frac1{5}V,\qquad
 D=\operatorname{diag}(- \frac{5}{4},0,\frac{3}{4},\frac{3}{4}).
\]
Dividing columns by these nonzero lengths gives $P^TP=I$, and the
eigenspace equations give $AP=PD$. Multiplying the latter by $P^T$
therefore verifies $A=PDP^T$.

\exercise{7.1.40}{Compute eigenvalues, then row-reduce each eigenspace}
The eigenvalues were computed from the characteristic polynomial. For
the eigenvectors we solve null-space systems by row reduction, not by
calling an eigenvector routine, as the exercise requests.
\step{Solve the eigenspace equations.}
$\lambda=-28$: row reduction gives \(x_{1}=- \frac{x_{5}}{4},\quad x_{2}=- \frac{x_{5}}{4},\quad x_{3}=- \frac{x_{5}}{4},\quad x_{4}=- \frac{x_{5}}{4}\).
$\lambda=8$: row reduction gives \(x_{1}=- x_{2} - x_{3},\quad x_{4}=0,\quad x_{5}=0\).
$\lambda=17$: row reduction gives \(x_{1}=x_{5},\quad x_{2}=x_{5},\quad x_{3}=x_{5},\quad x_{4}=x_{5}\).
$\lambda=32$: row reduction gives \(x_{1}=- \frac{x_{4}}{3},\quad x_{2}=- \frac{x_{4}}{3},\quad x_{3}=- \frac{x_{4}}{3},\quad x_{5}=0\).
\step{Normalize the orthogonal vectors and pair the same order with D.}
The columns of
\[
 V=\begin{bmatrix}1&1&1&1&1\\1&-1&1&1&1\\1&0&-2&1&1\\1&0&0&1&-3\\-4&0&0&1&0\end{bmatrix}
\]
are mutually orthogonal eigenvectors in the order
\[
 (\lambda_1,\ldots,\lambda_5)=(-28,8,8,17,32).
\]
Their squared lengths are $(20,2,6,5,12)$. Hence an explicit answer is
\[
 P=V\operatorname{diag}(\frac{\sqrt{5}}{10},\frac{\sqrt{2}}{2},\frac{\sqrt{6}}{6},\frac{\sqrt{5}}{5},\frac{\sqrt{3}}{6}),\qquad
 D=\operatorname{diag}(-28,8,8,17,32).
\]
Dividing columns by these nonzero lengths gives $P^TP=I$, and the
eigenspace equations give $AP=PD$. Multiplying the latter by $P^T$
therefore verifies $A=PDP^T$.
